% PDF 133--139；原书第 103--109 页。
\subsection{自动设计模板——离散搜索}
自动设计模板，在预先定义的搜索空间中用搜索或生成的方法，寻找效果最好的提示模板。这种方法实现起来更复杂，不过它一般会更充分地搜索参数空间，通常比手工提示工程表现更好。自动提示工程可分为离散搜索（discrete search）与连续搜索（continuous search）。

这些自动设计方法主要区别在于，使用离散的还是连续的词元或提示来引导语言模型。离散提示（discrete prompt）完全由自然语言词元或短语组成，连续提示（continuous prompt）则包含连续的、可调的参数。离散提示包括本节迄今所见的模板，其中连接输入 $x$ 与掩码输出 $z$ 的词元固定不变。连续提示中有不固定的词元，可在模型训练中调整。例如，可以用连续提示“[a1] [a2] [a3] [a4] [x] [a5] [z]”替换离散模板“the capital city of [x] is [z] .”，训练时微调其中的词元 $a_n$ 以改善结果。下面考察这些类别中的代表性方法及其应用前景。表~\ref{tab:ch03-discrete} 汇总了这些方法。

\subsubsection{提示挖掘}
提示挖掘（prompt mining）最早由 \refcite{ch03-jiang}{Jiang 等人（2020）}提出。它从大量文本语料中挖掘提示，依据是主语 $x$ 和宾语 $y$ 附近的词往往描述两者的关系。

\begin{keypoint}
仍以各国首都为例。在一个大型语料库中，Poland 和 Warsaw 在彼此附近出现时，通常暗示国家与首都之间的某种关系。如果收集许多具有同一种关系的主语—宾语对，例如更多国家与其首都，再抽取语料库中这些词对共同出现的句子，这些句子就可以作为生成该信息检索任务提示模板的依据。
\end{keypoint}

这种语料挖掘方法可以用两种办法生成提示。第一种称为中间词提示提取（middle-word prompt extraction）。在搜索语料中找到含有主语—宾语对的句子，抽取两者之间的文本词元，直接作为提示模板。仍以“波兰的首都是华沙”为例，目标是挖掘最能激发这条知识的提示。搜索语料库中包含这两个实体的句子，可以找到下面的例句。
\begin{quote}
\ttfamily
Warsaw is the capital city of Poland, and\\
has a population of 1.86 million people.
\end{quote}
只抽取主语—宾语对之间的词，得到“\texttt{is the capital city of}”。由此形成模板“\texttt{[z] is the capital city of [x]}”。

对小规模训练集得到的范围内所有主语—宾语对重复这个过程，搜索并抽取每一对的中间词提示。

\refcite{ch03-jiang}{Jiang 等人（2020）}的第二种挖掘方法使用更精细的语言学分析。具体而言，用句法分析从匹配的句子中抽取主语—宾语对之间最短的依存路径，将它作为模板。

例如，从句子“The capital of Poland is Warsaw.”抽取的中间词模板为“[x] is [z]”，这个模板显然太简单，无法完成任务。对同一个句子作依存分析，会得到下面的依存路径。
\[
\text{Poland}\xleftarrow{\mathrm{pobj}}\text{of}
\xleftarrow{\mathrm{prep}}\text{capital}\xleftarrow{\mathrm{nsubj}}\text{is}
\xleftarrow{\mathrm{attr}}\text{Warsaw}
\]
由此得到模板“capital of [x] is [z]”，看起来更适合查询一国首都。依存分析得到的模板也可能更容易激发目标知识，因为它们来自较为固定的语言规则，而大语言模型应已在预训练中学到这些规则（\refcite{ch03-jawahar}{Jawahar et al., 2019）}。

\subsubsection{提示改写}
提示改写（prompt paraphrasing）从已有提示出发，生成措辞尽可能多样的模板变体。仍以查询各国首都为例，可以写出若干略有区别的版本。
\begin{itemize}
\item \textbf{原始提示}\enspace “[z] is the capital city of [x] .”
\item \textbf{改写 1}\enspace “[z], the capital city of [x] .”
\item \textbf{改写 2}\enspace “[z] is the capital of [x] .”
\item \textbf{改写 3}\enspace “[x]’s capital city is [z] .”
\item \textbf{改写 $n$}\enspace “[x]’s capital, [z] .”
\end{itemize}
为自动完成这种改写，\refcite{ch03-mallinson}{Mallinson 等人（2017）}提出回译（back-translation）。先把原提示译成另一种语言，得到 $B$ 个候选译文。再把每个候选译文译回原来的语言，共得到 $B^2$ 个候选模板。按照往返翻译概率给模板排序，再从中筛选，概率计算如下。
\begin{equation}
P(t)=P_{\mathrm{forward}}(\bar{t}\mid\hat{t})\cdot P_{\mathrm{backward}}(t\mid\bar{t}),
\tag{3.9}
\end{equation}
其中 $\hat{t}$ 是原始提示，$\bar{t}$ 是原始提示 $\hat{t}$ 的译文，$t$ 是最终参与排序的候选提示。保留排名前 $T$ 的候选。之后还可进一步筛选和集成候选提示，使它们在目标自然语言处理任务中更有用。第 3.5 节会继续讨论提示筛选与集成。

\subsubsection{梯度引导搜索}
还可以用梯度引导搜索（gradient-directed search）设计提示模板。这个想法最早由 \refcite{ch03-wallace}{Wallace 等人（2019）}提出，他们当时研究对生成模型的对抗攻击（adversarial attack）。研究者设计了一种算法，把“触发词元”（trigger token）加到前缀提示末尾，紧挨答案空位之前，并迭代更新这些词元，使将错误回答填入空位时的损失最小。下面是一个问答（Question Answering，QA）提示的例子。
\begin{equation}
\text{“Question: [x] Context: [y] Answer: [T] [T] [T] }[z_{\mathrm{adv}}]\text{”},
\tag{3.10}
\end{equation}
其中 \texttt{[x]} 和 \texttt{[y]} 分别是问题与上下文，$[z_{\mathrm{adv}}]$ 是试图诱使模型产生的对抗输出，\texttt{[T]} 表示一串可变的触发词元。这些词元可迭代更新，以最小化某个语言模型对该序列的损失。更新使用基于 HotFlip 的梯度引导搜索（\refcite{ch03-ebrahimi}{Ebrahimi et al., 2018）}。这种做法诱使模型产生对抗回答。重要的是，研究者发现，在许多情形下，输入文本即使变化，优化出的触发词元序列仍然有效，会让模型针对许多不同输入生成同一个不合适的输出。下面是该研究使用 SQuAD 数据集中的一组问答所做的示例。

\begin{quote}
\textit{Question:} Why did he walk?

For exercise, Tesla walked between 8 to 10 miles per day. He squished his toes one hundred times for each foot every night, saying that it stimulated his brain cells. \textbf{why how because} \underline{to kill american people}
\end{quote}
三个词元“why how because”是梯度优化的产物。对许多不同输入，它们都会使 GPT-2 生成带下划线的对抗回答。

人们很快认识到，这种方法可以用于提示工程中的模板优化。在此基础上，\refcite{ch03-shin}{Shin 等人（2020）}提出了自动构造提示模板的方法 AutoPrompt。研究者用一组包含触发词元的初始模板，形式类似式（3.10），再在大量有标签的输入—输出样例上迭代，用梯度引导搜索优化词元。图~\ref{fig:ch03-autoprompt} 以情感分析为例说明了这一方法。送给语言模型的输入由三部分组成。
\begin{itemize}
\item \textbf{原始输入（$x_{\mathrm{inp}}$）。}对应图~\ref{fig:ch03-flow} 的输入 $x$。
\item \textbf{触发词元（$x_{\mathrm{trig}}$）。}通过梯度搜索学到的自然语言词元。学到多少词元取决于初始化梯度搜索时设定的词元数，这个数量可以视作超参数。
\item \textbf{答案空位。}在图~\ref{fig:ch03-autoprompt} 中记作 \texttt{[P]} 或 \texttt{[MASK]}，对应图~\ref{fig:ch03-flow} 示例中的 \texttt{[z]}。
\end{itemize}

\sourcefigure{fig03_11.pdf}{用 AutoPrompt 考察掩码语言模型能否进行情感分析。输入 $x_{\mathrm{inp}}$ 被放入自然语言提示 $x_{\mathrm{prompt}}$，其中有一个待预测的 \texttt{[MASK]} 词元。模板 $\lambda$ 将原始输入与学到的触发词元 $x_{\mathrm{trig}}$ 组合起来。触发词元在各个输入间共享，由上述梯度搜索过程学得。对各类别的自动选择的标签词元集合，将掩码语言模型预测的 $p(\texttt{[MASK]}\mid x_{\mathrm{prompt}})$ 求和，得到类别 $y$ 的概率。}{fig:ch03-autoprompt}

按给定提示模板的结构组合各部件，得到送入语言模型的优化输入 $x_{\mathrm{prompt}}$。将自动选择的一组输出词元的概率相加，就能确定标签类别。在本例中，正面类别对应 Cris、marvelous 和 philanthrope，负面类别对应 worse、incompetence 和 Worse。正面标签的累计概率高于负面标签，因此分类结果为正面。

\begin{keypoint}
优化得到的词元在人看来可能不直观。在同一个 LAMA T-REx 实体关系子集基准上，\refcite{ch03-shin}{Shin 等人（2020）}报告称，以手工模板为参照，AutoPrompt 比 \refcite{ch03-jiang}{Jiang 等人（2020）}评估的 Top-1 改写提示多提高了 9 个百分点的准确率。使用 BERT 和 RoBERTa 变体时，在答案映射、自然语言推断、事实检索和实体关系抽取等任务上，AutoPrompt 的表现比手工提示高 10\% 至 20\%。经过优化的提示甚至可以超过微调模型，尤其是在数据较少、只有少数有标签样例的场景中。
\end{keypoint}

\subsubsection{提示生成}
上述离散提示搜索方法都使用掩码语言模型，预测单个词元。受 \refcite{ch03-radford-2019}{Radford 等人（2019）}、\refcite{ch03-brown}{Brown 等人（2020）}及其他研究展示的上下文学习能力启发，\refcite{ch03-gao}{Gao 等人（2021）}提出了语言模型的改进少样本微调（Better Few-shot Fine-tuning of Language Models，LM-BFF）。他们用大型文本到文本模型 T5（\refcite{ch03-raffel}{Raffel et al., 2020）}取代前面例子中的掩码语言模型。其方法是以少量训练样例组成少样本提示，引导 T5 搜索合适的提示。

图~\ref{fig:ch03-lmbff} 展示 \refcite{ch03-gao}{Gao 等人（2021）}用 T5 生成提示模板的过程。先从小型训练集 $D_{\mathrm{train}}$ 抽取若干样例，送入 T5，提示它构造模板 $\mathcal{T}$，不必明确指定模板应包含多少个词元。T5 的输入经过专门设计，以引导模型生成提示。图中将 $D_{\mathrm{train}}$ 的样例按类别分组，组织成适合 T5 的输入，再生成一个模板或一组模板 $\mathcal{T}$，使下式最大。
\begin{equation}
\sum_{(x_{\mathrm{in}},y)\in D_{\mathrm{train}}}
\log P_{\mathrm{T5}}\bigl(\mathcal{T}\mid\mathcal{T}(x_{\mathrm{in}},y)\bigr),
\tag{3.11}
\end{equation}
其中 $P_{\mathrm{T5}}$ 是 T5 语言模型的输出概率分布。

\sourcefigure{fig03_12.pdf}{提示生成过程。把训练集 $D_{\mathrm{train}}$ 的输入样例按类别划分，构造为适合 T5 的输入，通过解码得到一个或少量模板，使式（3.11）的求和最大。图中先生成候选模板，再经微调和评估选出最佳模板。}{fig:ch03-lmbff}

在用这些构造出的输入使 T5 生成一组模板后，解码生成的模板。在 $D_{\mathrm{train}}$ 上对语言模型 $\mathcal{L}$ 作微调，再用开发集 $D_{\mathrm{dev}}$ 评估，选出最佳模板。\refcite{ch03-gao}{Gao 等人（2021）}表明，这种提示模板生成方法与在给定输入旁附上语义相似的示范相结合，表现显著优于手工设计的提示。此外，用 T5 配合手工生成或自动生成的提示，表现优于标准微调方案，说明了基于提示的学习在自然语言处理中的作用。

在 \refcite{ch03-gao}{Gao 等人（2021）}的工作基础上，\refcite{ch03-ben-david}{Ben-David 等人（2022）}提出即时任意领域适配的提示学习算法（Prompt learning algorithm for on-the-fly Any-Domain Adaptation，PADA）。该方法也利用 T5 生成文本到文本提示的能力，专门生成可以表示多个源领域、供人阅读的提示。针对预测分布外数据（out-of-distribution data）这一常见难题，PADA 将多个特定领域映射到同一个语义空间，以提高泛化能力。\refcite{ch03-ben-david}{Ben-David 等人（2022）}报告称，在文本分类和序列标注任务上，该方法相对于表现良好的对照方法取得了明显提升。

% 按原书第 109 页重排表 3.3，保留横置阅读方向。
\begin{table}[p]
\centering
\makebox[\linewidth][c]{\rotatebox{90}{%
\begin{minipage}{.94\textheight}
\caption{离散提示的自动设计方法。}
\label{tab:ch03-discrete}
\footnotesize
\setlength{\tabcolsep}{3pt}
\renewcommand{\arraystretch}{1.14}
\begin{tabularx}{\linewidth}{@{}>{\raggedright\arraybackslash}p{.10\linewidth}
  >{\raggedright\arraybackslash}p{.155\linewidth}
  >{\raggedright\arraybackslash}X
  >{\raggedright\arraybackslash}p{.174\linewidth}
  >{\raggedright\arraybackslash}p{.182\linewidth}@{}}
\toprule
方法 & 概述 & 示例 & 优点 & 缺点\\
\midrule
提示挖掘 &
在大型语料中搜索包含给定查询—答案对的句子，并用它们构造提示模板。 &
\textbf{测试对}\enspace Poland :: Warsaw\par
\textbf{挖掘到的句子}\enspace “Warsaw is the capital city of Poland, and has a population 1.86 million people”（见维基百科）\par
\textbf{所得模板}\enspace “[z] is the capital city of [x] .” &
程序实现简单。\par 可用更复杂的语言学抽取过程扩展。 &
所得模板受语料中已有句子的领域限制。\par 难以保证优化效果。\\
\midrule
提示改写 &
让种子提示经过多轮翻译，生成许多略有差异的变体，再选出表现最好的模板。 &
\textbf{种子提示}\enspace [z] is the capital city of [x] .\par
\textbf{回译变体}\par
“[z], the capital city of [x] .”\par
“[x]’s capital city is [z] .”\par
“[x]’s capital, [z] .” &
程序实现简单。\par 测试的提示多样，有助于优化。 &
回答覆盖范围较窄。\par 难以保证优化效果。\\
\midrule
梯度搜索 &
将一组可变的触发词元（本例为 [T]）和输入—预测对组合成模板，在训练中优化，使预测结果最好。 &
\textbf{评论／情感对}\enspace “a real joy” :: positive\par
\textbf{初始提示}\enspace “a real joy [T][T][T] [T][T] positive”\par
\textbf{梯度优化后的提示}\enspace “a real joy atmosphere alot dialogue Clone totally positive”\par
\textbf{优化后的模板}\enspace [x] atmosphere alot dialogue Clone totally [z] &
可以找到高度优化的输入词元。\par 不依赖该领域已有的句子语料。 &
计算开销大，程序实现复杂。\par 产生的模板不直观。\par 输出模板受人类语言的嵌入表示约束。\\
\midrule
提示生成 &
编码器—解码器模型（如 T5）预测种子模板中的词元。种子模板由训练查询—回答对构造。再把预测结果转成模板并评估质量。 &
\textbf{影评情感的种子模板}\par
A pleasure to watch. <X> great <Y>\par
No reason to watch. <X> terrible <Y>\par
\textbf{T5 填充后的模板}\par
A pleasure to watch. This is great .\par
No reason to watch. A terrible one.\par
\textbf{所得待测模板}\par
[x] This is [z].\par
[x] A [z] one. &
测试的提示多样，有助于优化。\par 大语言模型生成的模板本身可能更适合大语言模型使用。 &
计算开销大，程序实现复杂。\par 种子模板需要人工输入。\par 难以保证优化效果。\\
\bottomrule
\end{tabularx}
\end{minipage}}}
\end{table}

